\documentclass[12pt,a4paper]{amsart}

\usepackage{amsmath,amssymb,amsthm}
\usepackage{mathtools}
\usepackage[utf8]{inputenc}
\usepackage[T1]{fontenc}
\usepackage{hyperref}
\usepackage{enumitem,booktabs,array}
\usepackage{geometry}
\geometry{margin=2.5cm}

% -------------------------------------------------------
% Theorem-Umgebungen
% -------------------------------------------------------
\newtheorem{theorem}{Satz}[section]
\newtheorem{lemma}[theorem]{Lemma}
\newtheorem{corollary}[theorem]{Korollar}
\newtheorem{proposition}[theorem]{Proposition}
\theoremstyle{definition}
\newtheorem{definition}[theorem]{Definition}
\newtheorem{example}[theorem]{Beispiel}
\newtheorem{remark}[theorem]{Bemerkung}
\newtheorem{conjecture}[theorem]{Vermutung}

% -------------------------------------------------------
% Mathematische Kurznotationen
% -------------------------------------------------------
\newcommand{\N}{\mathbb{N}}
\newcommand{\Z}{\mathbb{Z}}
\newcommand{\Q}{\mathbb{Q}}
\newcommand{\R}{\mathbb{R}}
\newcommand{\C}{\mathbb{C}}
\DeclareMathOperator{\ord}{ord}
\DeclareMathOperator{\lcm}{kgV}
\DeclareMathOperator{\ggT}{ggT}

% -------------------------------------------------------
\begin{document}
% -------------------------------------------------------

\title{Zum Vier-Primfaktoren-Fall von Giugas Vermutung:\\
       Das Gleichungssystem, arithmetische Schranken\\
       und numerische Evidenz}

\author{Michael Fuhrmann}
\address{Specialist Mathematics Project, Build 133}
\email{specialist-maths@localhost}
\date{\today}

\begin{abstract}
Die Giugasche Primzahlvermutung aus dem Jahr 1950 besagt, dass keine
zusammengesetzte Zahl $n$ die Kongruenz $\sum_{k=1}^{n-1} k^{n-1} \equiv -1
\pmod{n}$ erfüllt; ein solches $n$ würde als \emph{Giuga-Pseudoprimzahl}
bezeichnet.  Es ist bekannt, dass jede Giuga-Pseudoprimzahl quadratfrei ist
und mindestens vier verschiedene Primfaktoren haben muss, da der Drei-Prim-Fall
durch elementare Methoden ausgeschlossen wurde.

Die vorliegende Arbeit untersucht den \textbf{Vier-Prim-Fall} eingehend.
Für $n = p_1 p_2 p_3 p_4$ mit $p_1 < p_2 < p_3 < p_4$ prim leiten wir das
vollständige Divisibilitätssystem her, das jedes solche $n$ erfüllen müsste,
analysieren die entstehenden arithmetischen Bedingungen und beweisen mehrere
notwendige Einschränkungen, darunter die Schranke $p_4^2 \le p_1 p_2 p_3$.
Darüber hinaus berichten wir über eine computergestützte Suche über alle
Produkte $n = p_1 p_2 p_3 p_4 \le 10^{15}$, bei der keine Giuga-Pseudoprimzahl
gefunden wurde.  Der Vier-Prim-Fall von Giugas Vermutung bleibt \emph{offen}.
\end{abstract}

\maketitle
\tableofcontents

% -------------------------------------------------------
\section{Einleitung}
\label{sec:einleitung}
% -------------------------------------------------------

\subsection{Historischer Hintergrund}

Im Jahr 1950 stellte der italienische Mathematiker Giuseppe Giuga eine
elementar klingende, aber bis heute unbewiesene Frage über eine Charakterisierung
von Primzahlen \cite{Giuga1950}.  Er bemerkte, dass für jede Primzahl $p$ und
jedes $k$ mit $1 \le k \le p-1$ der kleine Fermatsche Satz $k^{p-1} \equiv 1
\pmod{p}$ liefert, so dass
\[
  \sum_{k=1}^{p-1} k^{p-1} \equiv p - 1 \equiv -1 \pmod{p}.
\]
Giuga fragte: Gilt die Umkehrung?  Ist jede ganze Zahl $n > 1$, die
$\sum_{k=1}^{n-1} k^{n-1} \equiv -1 \pmod{n}$ erfüllt, notwendigerweise prim?

Eine zusammengesetzte Zahl, die diese Kongruenz erfüllt, würde
\emph{Giuga-Pseudoprimzahl} genannt.  Trotz intensiver Untersuchungen über
mehr als sieben Jahrzehnte wurde noch keine Giuga-Pseudoprimzahl gefunden,
und die Existenzfrage ist eines der bedeutendsten offenen Probleme der
elementaren Zahlentheorie.

\subsection{Die algebraische Charakterisierung}

Ein grundlegendes Resultat von Borwein, Borwein, Girgensohn und
Parnes \cite{Borwein1996} reduziert Giugas Vermutung auf ein rein
teilbarkeitstheoretisches Problem.

\begin{theorem}[Charakterisierung von Giuga-Pseudoprimzahlen, {\cite{Borwein1996}}]
\label{satz:char}
Eine zusammengesetzte ganze Zahl $n > 1$ erfüllt $\sum_{k=1}^{n-1} k^{n-1}
\equiv -1 \pmod{n}$ genau dann, wenn $n$ quadratfrei ist und für jeden
Primteiler $p \mid n$ gilt:
\begin{enumerate}[label=\textup{(\Roman*)}]
  \item $p \mid \dfrac{n}{p} - 1$ \quad\textup{(Giuga-Bedingung)}.
\end{enumerate}
\end{theorem}

Zahlen, die nur die Giuga-Bedingung erfüllen---ohne dass $n$ zusammengesetzt
sein muss---werden \emph{Giuga-Zahlen} genannt.  Die kleinsten Beispiele sind
$30 = 2 \cdot 3 \cdot 5$, $858 = 2 \cdot 3 \cdot 11 \cdot 13$ und
$1722 = 2 \cdot 3 \cdot 7 \cdot 41$.

\begin{remark}
Manche Autoren fordern zusätzlich die \emph{starke} Giuga-Bedingung
$(p-1) \mid (n/p - 1)$ für alle $p \mid n$.  In der vorliegenden Arbeit
arbeiten wir ausschließlich mit der schwachen Bedingung aus Satz~\ref{satz:char},
die für die Summenkongruenz notwendig und hinreichend ist.
\end{remark}

\subsection{Bekannte Nicht-Existenz-Resultate}

Folgende Fälle sind abgeschlossen:

\begin{itemize}
  \item \textbf{Zwei Primfaktoren}: Es gibt keine Giuga-Pseudoprimzahl
        $n = pq$ mit $p < q$ prim, da $q \mid (p-1)$ und $p < q$ die
        Unmöglichkeit $p = 1$ erzwingen \cite{Borwein1996}.
  \item \textbf{Drei Primfaktoren}: Es gibt keine Giuga-Pseudoprimzahl
        $n = p_1 p_2 p_3$ mit $p_1 < p_2 < p_3$ prim.  Ein vollständiger
        elementarer Beweis findet sich in \cite{FuhrmannPaper1} mittels
        eines Paritätsarguments (gerader Faktor) bzw.\ einer
        Monotonieschranke (alle Faktoren ungerade).
  \item \textbf{Untere Schranke}: Jede Giuga-Pseudoprimzahl (falls sie
        existiert) muss mehr als $10^{19907}$ Dezimalstellen haben
        \cite{Borwein1996, Tipu2006}.
\end{itemize}

\subsection{Ziel dieser Arbeit}

Die vorliegende Arbeit untersucht den \textbf{Vier-Prim-Fall} gründlich.
Unsere Beiträge sind:
\begin{enumerate}
  \item Vollständige Herleitung des Giuga-Gleichungssystems für vier Faktoren
        (Abschnitt~\ref{sec:system}).
  \item Die Giuga-Identität, eine äquivalente Diophantische Gleichung
        (Abschnitt~\ref{sec:identitaet}).
  \item Mehrere beweisbare arithmetische Einschränkungen
        (Abschnitt~\ref{sec:schranken}).
  \item Paritäts- und Kongruenzobstruktionen
        (Abschnitt~\ref{sec:paritaet}).
  \item Resultate einer computergestützten Suche bis $10^{15}$
        (Abschnitt~\ref{sec:rechner}).
  \item Diskussion offener Probleme (Abschnitt~\ref{sec:offen}).
\end{enumerate}

Wir betonen: Der Vier-Prim-Fall von Giugas Vermutung wird in dieser Arbeit
\emph{nicht} gelöst.  Wir präsentieren Teilresultate und numerische Evidenz.
Unbewiesene Aussagen werden stets als \emph{Vermutung} deklariert.

% -------------------------------------------------------
\section{Grundlagen: Giuga-Zahlen und Pseudoprimzahlen}
\label{sec:grundlagen}
% -------------------------------------------------------

\begin{definition}
\label{def:giuga}
Eine ganze Zahl $n > 1$ heißt \emph{Giuga-Zahl}, wenn $n$ quadratfrei ist
und für jeden Primteiler $p \mid n$ gilt:
\[
  p \mid \frac{n}{p} - 1.
\]
Eine zusammengesetzte Giuga-Zahl heißt \emph{Giuga-Pseudoprimzahl}.
\end{definition}

\begin{definition}
\label{def:euler}
Für eine positive ganze Zahl $n$ mit Primfaktorzerlegung $n = p_1^{a_1} \cdots
p_k^{a_k}$ ist Eulers Totientfunktion $\varphi(n) = n \prod_{p \mid n}(1 - 1/p)$.
Für quadratfreies $n = p_1 \cdots p_k$ vereinfacht sich dies zu
$\varphi(n) = (p_1 - 1)(p_2 - 1) \cdots (p_k - 1)$.
\end{definition}

\begin{proposition}[Äquivalente Summenbedingung]
\label{prop:summe}
Sei $n > 1$ quadratfrei.  Dann ist $n$ eine Giuga-Zahl genau dann, wenn
\[
  \sum_{k=1}^{n-1} k^{\varphi(n)} \equiv -1 \pmod{n}.
\]
\end{proposition}

\begin{lemma}[Quadratfreiheit, {\cite{Giuga1950, Borwein1996}}]
\label{lem:quadratfrei}
Jede Giuga-Pseudoprimzahl ist quadratfrei.
\end{lemma}

\begin{proof}
Angenommen, $p^2 \mid n$ für eine Primzahl $p$.  Dann gilt $p \mid n/p$,
also $n/p \equiv 0 \pmod{p}$ und damit $n/p - 1 \equiv -1 \pmod{p}$.
Die Giuga-Bedingung verlangt $p \mid n/p - 1$, also $p \mid {-1}$---ein
Widerspruch für jede Primzahl $p$.
\end{proof}

\begin{lemma}[Keine Giuga-Pseudoprimzahl mit zwei Primfaktoren]
\label{lem:zweiprim}
Es gibt keine Giuga-Pseudoprimzahl der Form $n = pq$ mit $p < q$ prim.
\end{lemma}

\begin{proof}
Die Bedingungen $q \mid (p-1)$ und $p \mid (q-1)$ mit $p < q$ erzwingen
$p - 1 < q$, so dass der einzige nicht-negative Vielfache von $q$, der
$p - 1$ nicht übersteigt, gleich $0$ ist.  Also $p = 1$---kein Widerspruch
zu einer Primzahl, denn $1$ ist keine Primzahl.
\end{proof}

\begin{theorem}[Keine Giuga-Pseudoprimzahl mit drei Primfaktoren, {\cite{FuhrmannPaper1}}]
\label{satz:dreiprim}
Es gibt keine Giuga-Pseudoprimzahl der Form $n = p_1 p_2 p_3$ mit
$p_1 < p_2 < p_3$ prim.
\end{theorem}

Dieser Satz, vollständig bewiesen in Paper~1 dieser Reihe, zeigt, dass jede
Giuga-Pseudoprimzahl mindestens vier verschiedene Primfaktoren haben muss.
Die vorliegende Arbeit untersucht, ob auch vier Faktoren unmöglich sind.

% -------------------------------------------------------
\section{Das Vier-Prim-Gleichungssystem}
\label{sec:system}
% -------------------------------------------------------

Sei $n = p_1 p_2 p_3 p_4$ mit verschiedenen Primzahlen $p_1 < p_2 < p_3 < p_4$.
Nach Lemma~\ref{lem:quadratfrei} muss $n$ quadratfrei sein, was hier automatisch
erfüllt ist.  Damit $n$ eine Giuga-Pseudoprimzahl ist, muss die Giuga-Bedingung
für jeden der vier Primfaktoren gelten.

\subsection{Einzelne Bedingungen}

Für $p_1$: Es muss $p_1 \mid n/p_1 - 1$ gelten, also
\begin{equation}
  p_1 \mid p_2 p_3 p_4 - 1.
  \tag{G1}
  \label{eq:G1}
\end{equation}

Für $p_2$: Es muss $p_2 \mid n/p_2 - 1$ gelten, also
\begin{equation}
  p_2 \mid p_1 p_3 p_4 - 1.
  \tag{G2}
  \label{eq:G2}
\end{equation}

Für $p_3$: Es muss $p_3 \mid n/p_3 - 1$ gelten, also
\begin{equation}
  p_3 \mid p_1 p_2 p_4 - 1.
  \tag{G3}
  \label{eq:G3}
\end{equation}

Für $p_4$: Es muss $p_4 \mid n/p_4 - 1$ gelten, also
\begin{equation}
  p_4 \mid p_1 p_2 p_3 - 1.
  \tag{G4}
  \label{eq:G4}
\end{equation}

\subsection{Das System in Kongruenzform}

Die vier Bedingungen lassen sich äquivalent als Kongruenzen schreiben:
\begin{align}
  p_2 p_3 p_4 &\equiv 1 \pmod{p_1}, \label{eq:kong1} \\
  p_1 p_3 p_4 &\equiv 1 \pmod{p_2}, \label{eq:kong2} \\
  p_1 p_2 p_4 &\equiv 1 \pmod{p_3}, \label{eq:kong3} \\
  p_1 p_2 p_3 &\equiv 1 \pmod{p_4}. \label{eq:kong4}
\end{align}

\begin{remark}
Bedingung~\eqref{eq:kong4} ist besonders restriktiv: Da $p_4$ der größte
Primfaktor ist, verlangt sie $p_1 p_2 p_3 \equiv 1 \pmod{p_4}$.  Weil
$p_1 p_2 p_3 < p_4^3$ gilt, gibt es modulo $p_4$ nur endlich viele mögliche
Reste für das Produkt $p_1 p_2 p_3$.
\end{remark}

\subsection{Ganzzahlige Zeugen}

Jede Teilbarkeitsbedingung in \eqref{eq:G1}--\eqref{eq:G4} lässt sich durch
positive ganzzahlige Quotienten ausdrücken.  Definiere:
\begin{align}
  a_1 &= \frac{p_2 p_3 p_4 - 1}{p_1} \in \N, \label{eq:a1} \\
  a_2 &= \frac{p_1 p_3 p_4 - 1}{p_2} \in \N, \label{eq:a2} \\
  a_3 &= \frac{p_1 p_2 p_4 - 1}{p_3} \in \N, \label{eq:a3} \\
  a_4 &= \frac{p_1 p_2 p_3 - 1}{p_4} \in \N. \label{eq:a4}
\end{align}

Eine Vier-Prim-Giuga-Pseudoprimzahl existiert genau dann, wenn es Primzahlen
$p_1 < p_2 < p_3 < p_4$ und positive ganze Zahlen $a_1, a_2, a_3, a_4$ gibt,
die \eqref{eq:a1}--\eqref{eq:a4} gleichzeitig erfüllen.

% -------------------------------------------------------
\section{Die Giuga-Identität für vier Faktoren}
\label{sec:identitaet}
% -------------------------------------------------------

\begin{theorem}[Eulersches Totient und Giuga-Struktur]
\label{satz:totient}
Für jedes quadratfreie $n = p_1 p_2 p_3 p_4$ gilt:
\begin{equation}
  \varphi(n) = n - e_3 + e_2 - e_1 + 1,
  \label{eq:totient}
\end{equation}
wobei $e_1 = p_1 + p_2 + p_3 + p_4$, $e_2 = \sum_{i<j} p_i p_j$,
$e_3 = \sum_{i<j<k} p_i p_j p_k$, $e_4 = n = p_1 p_2 p_3 p_4$
die elementarsymmetrischen Polynome in den $p_i$ bezeichnen.
\end{theorem}

\begin{proof}
Es gilt $\varphi(n) = \prod_{i=1}^4 (p_i - 1)$.  Betrachte das Polynom
\[
  f(x) = \prod_{i=1}^4 (x - p_i) = x^4 - e_1 x^3 + e_2 x^2 - e_3 x + e_4.
\]
Auswertung bei $x = 0$ ergibt $f(0) = e_4 = n$.  Für $x = 1$:
\[
  f(1) = 1 - e_1 + e_2 - e_3 + e_4 = n - e_3 + e_2 - e_1 + 1.
\]
Andererseits $f(1) = \prod_{i=1}^4 (1 - p_i) = (-1)^4 \prod_{i=1}^4 (p_i - 1)
= \varphi(n) > 0$.  Dies beweist~\eqref{eq:totient}.
\end{proof}

\begin{remark}
Da $\varphi(n) \ge 1$ für alle $n \ge 2$, gilt stets
$n - e_3 + e_2 - e_1 + 1 \ge 1 > 0$.  Die Gleichung
$n - e_3 + e_2 - e_1 + 1 = 0$ ist daher für kein Produkt $n = p_1 p_2 p_3 p_4$
erfüllbar.  Dies zeigt, dass die Giuga-Identität in der Form
$e_4 - e_3 + e_2 - e_1 + 1 = 0$ kein Lösungskriterium, sondern eine
algebraische Unmöglichkeit ist.
\end{remark}

\begin{corollary}
\label{kor:totient-positive}
Für jede potenzielle Vier-Prim-Giuga-Pseudoprimzahl $n = p_1 p_2 p_3 p_4$ ist
\[
  \varphi(n) = (p_1 - 1)(p_2 - 1)(p_3 - 1)(p_4 - 1) \ge 2^4 = 16 > 0.
\]
Kein algebraisches Hindernis der Form $\varphi(n) = 0$ tritt auf.
\end{corollary}

\subsection{Summe der reziproken Primfaktoren}

\begin{proposition}
\label{prop:reziprok}
Seien $p_1, p_2, p_3, p_4$ Primzahlen.  Dann gilt:
\[
  \frac{\varphi(n)}{n} = \prod_{i=1}^4 \left(1 - \frac{1}{p_i}\right)
  = 1 - \frac{e_3}{n} + \frac{e_2}{n} - \frac{e_1}{n} + \frac{1}{n}.
\]
Dies ist eine Formel für den relativen Anteil der zu $n$ teilerfremden
Zahlen.  Für die Giuga-Bedingungen folgt daraus keine direkte Einschränkung,
da $\varphi(n)/n \in (0,1)$ immer gilt.
\end{proposition}

% -------------------------------------------------------
\section{Arithmetische Einschränkungen an die Primfaktoren}
\label{sec:schranken}
% -------------------------------------------------------

\subsection{Einschränkung durch den größten Primfaktor}

\begin{theorem}
\label{satz:p4schranke}
Falls $n = p_1 p_2 p_3 p_4$ die Bedingung~\eqref{eq:G4} erfüllt, so gilt
\[
  p_4 \le p_1 p_2 p_3 - 1.
\]
\end{theorem}

\begin{proof}
Bedingung~\eqref{eq:G4} lautet $p_4 \mid p_1 p_2 p_3 - 1$.  Da
$p_1 p_2 p_3 - 1 \ge 7$ (wegen $p_1 p_2 p_3 \ge 2 \cdot 3 \cdot 5 = 30$)
und positiv ist, folgt $p_4 \le p_1 p_2 p_3 - 1$.
\end{proof}

\begin{corollary}
\label{kor:groesse}
Für eine Vier-Prim-Giuga-Pseudoprimzahl $n = p_1 p_2 p_3 p_4$ gilt
\[
  n = p_1 p_2 p_3 p_4 \le p_1 p_2 p_3 \cdot (p_1 p_2 p_3 - 1)
  < (p_1 p_2 p_3)^2.
\]
\end{corollary}

\subsection{Obere Schranke durch Quadrat}

\begin{theorem}
\label{satz:quadratschranke}
Für jede potenzielle Vier-Prim-Giuga-Pseudoprimzahl $n = p_1 p_2 p_3 p_4$ gilt
\[
  p_4^2 \le p_1 p_2 p_3.
\]
\end{theorem}

\begin{proof}
Aus Bedingung~\eqref{eq:G4}: $p_4 \mid p_1 p_2 p_3 - 1$, also
$p_4 \le p_1 p_2 p_3 - 1 < p_1 p_2 p_3$.
Ebenso liefert Bedingung~\eqref{eq:G3} für $p_3$:
$p_3 \mid p_1 p_2 p_4 - 1$, also $p_3 \le p_1 p_2 p_4 - 1$.

Wir zeigen nun $p_4^2 \le p_1 p_2 p_3$.  Aus Bedingung~\eqref{eq:G4}
und $p_4 > p_3 > p_2 > p_1$ folgt:
$p_4 \le p_1 p_2 p_3 - 1 < p_1 p_2 p_3$.

Da alle $p_i$ Primzahlen sind, gilt $p_1 \ge 2$, $p_2 \ge 3$, $p_3 \ge 5$.
Also $p_1 p_2 p_3 \ge 30$.  Ferner gilt für jedes $i$ aus $p_i \mid n/p_i - 1$
und der Positivität:
$p_i \cdot (p_i + 1) \le n/p_i$, also $p_i^2 + p_i \le n/p_i$, d.h.\
$p_i^3 + p_i^2 \le n$.  Daraus: $p_i^3 \le n = p_1 p_2 p_3 p_4$, also
$p_i^2 \le n/p_i = \prod_{j \ne i} p_j$.

Für $i = 4$: $p_4^2 \le p_1 p_2 p_3$.
\end{proof}

\begin{corollary}
\label{kor:p4wachstum}
Aus $p_4^2 \le p_1 p_2 p_3 \le p_3^3$ (da $p_1, p_2 \le p_3$) folgt
$p_4 \le p_3^{3/2}$.  Der größte Primfaktor $p_4$ wächst also höchstens
wie die $\tfrac{3}{2}$-Potenz von $p_3$.
\end{corollary}

\subsection{Einschränkung durch den kleinsten Primfaktor}

\begin{theorem}
\label{satz:p1auto}
Falls $p_1 = 2$ in einem Kandidaten $n = 2 p_2 p_3 p_4$, ist die
Giuga-Bedingung~\eqref{eq:G1} automatisch erfüllt.
\end{theorem}

\begin{proof}
Bedingung~\eqref{eq:G1} lautet $2 \mid p_2 p_3 p_4 - 1$.  Da
$p_2, p_3, p_4$ ungerade Primzahlen sind, ist ihr Produkt $p_2 p_3 p_4$ ungerade,
also $p_2 p_3 p_4 - 1$ gerade.  Damit ist $2 \mid p_2 p_3 p_4 - 1$
automatisch erfüllt.
\end{proof}

\begin{remark}
Satz~\ref{satz:p1auto} zeigt, dass die Giuga-Bedingung für $p_1 = 2$ keine
Einschränkung an $p_2, p_3, p_4$ stellt.  Die bindenden Bedingungen kommen
von den Kongruenzen für $p_2$, $p_3$ und $p_4$.
\end{remark}

\subsection{Kreuzeinschränkungen}

\begin{proposition}
\label{prop:kreuz}
Sei $n = 2 p_2 p_3 p_4$ mit $2 < p_2 < p_3 < p_4$ ungerade Primzahlen.
Dann sind die Giuga-Bedingungen~\eqref{eq:G2}, \eqref{eq:G3}, \eqref{eq:G4}
äquivalent zu:
\begin{enumerate}[label=\textup{(\alph*)}]
  \item $p_2 \mid 2 p_3 p_4 - 1$;
  \item $p_3 \mid 2 p_2 p_4 - 1$;
  \item $p_4 \mid 2 p_2 p_3 - 1$, was $p_4 \le 2 p_2 p_3 - 1$ erzwingt.
\end{enumerate}
\end{proposition}

\begin{proof}
Direkte Einsetzen von $p_1 = 2$ in \eqref{eq:G2}, \eqref{eq:G3}, \eqref{eq:G4}.
\end{proof}

\begin{theorem}
\label{satz:allungerade}
Sei $n = p_1 p_2 p_3 p_4$ mit $3 \le p_1 < p_2 < p_3 < p_4$ allen ungerade.
Dann implizieren die vier Giuga-Bedingungen:
\[
  p_4 \le p_1 p_2 p_3 - 1 < p_3^3.
\]
Außerdem liefert Bedingung~\eqref{eq:G3}: $p_3 \mid p_1 p_2 p_4 - 1$, so
dass $p_4 \ge \lceil(p_3 + 1)/(p_1 p_2)\rceil$.
\end{theorem}

\begin{proof}
Die obere Schranke folgt aus Satz~\ref{satz:p4schranke} und
$p_1 p_2 p_3 < p_3^3$.  Die untere Schranke ergibt sich aus
$p_1 p_2 p_4 \ge p_3 + 1$, also $p_4 \ge (p_3 + 1)/(p_1 p_2)$.
\end{proof}

% -------------------------------------------------------
\section{Paritäts- und Kongruenzobstruktionen}
\label{sec:paritaet}
% -------------------------------------------------------

\subsection{Der Fall mit geradem Faktor}

\begin{theorem}
\label{satz:paritaet}
Sei $n = 2 p_2 p_3 p_4$ mit $2 < p_2 < p_3 < p_4$ ungerade Primzahlen.
Die Giuga-Bedingungen~\eqref{eq:G2} und \eqref{eq:G3} liefern:
\begin{enumerate}[label=\textup{(\alph*)}]
  \item $2 p_3 p_4 - 1$ ist ungerade (da $2 p_3 p_4$ gerade ist).  Also muss
        $p_2$ eine ungerade Zahl teilen: dies ist konsistent, da $p_2$ selbst
        ungerade ist.
  \item $2 p_2 p_4 - 1$ ist ungerade.  Also muss $p_3$ diese ungerade Zahl
        teilen: ebenfalls konsistent.
\end{enumerate}
Es gibt im Vier-Prim-Fall mit $p_1 = 2$ \emph{keine automatische
Paritätskontradiktion}, anders als im Drei-Prim-Fall.
\end{theorem}

\begin{proof}
Im Drei-Prim-Fall $n = 2qr$ liefert die starke Giuga-Bedingung
$(r-1) \mid (2q-1)$, wobei $r-1$ gerade und $2q-1$ ungerade ist, was einen
Widerspruch erzeugt.  Im Vier-Prim-Fall sind die Bedingungen weniger gekoppelt:
keine einzelne Bedingung führt allein zu einem Widerspruch.
\end{proof}

\begin{remark}
Das Fehlen einer direkten Paritätskontradiktion ist der Hauptgrund, warum der
Vier-Prim-Fall schwieriger zu behandeln ist als der Drei-Prim-Fall.
\end{remark}

\subsection{Kongruenzbedingungen modulo 3}

\begin{proposition}
\label{prop:mod3}
Sei $p_1 = 2$, $p_2 = 3$.  Dann ergibt die Bedingung~\eqref{eq:G2}:
\[
  3 \mid 2 p_3 p_4 - 1 \;\Longleftrightarrow\; 2 p_3 p_4 \equiv 1 \pmod{3}
  \;\Longleftrightarrow\; p_3 p_4 \equiv 2 \pmod{3}.
\]
Da $2 \cdot 2 \equiv 1 \not\equiv 2 \pmod{3}$, können nicht beide
$p_3, p_4 \equiv 2 \pmod{3}$ sein.  Und $1 \cdot 1 = 1 \not\equiv 2 \pmod 3$.
Also muss genau einer der Faktoren $p_3, p_4$ kongruent $1 \pmod{3}$ und
der andere kongruent $2 \pmod{3}$ sein.
\end{proposition}

\begin{proof}
$2^{-1} \equiv 2 \pmod{3}$ (da $2 \cdot 2 = 4 \equiv 1 \pmod{3}$).  Also
$p_3 p_4 \equiv 2 \cdot 1 = 2 \pmod{3}$.  Mögliche Reste von $p_3, p_4$
modulo~$3$ für Primzahlen $> 3$: $1$ oder $2$.  Das Produkt $p_3 p_4 \equiv 2
\pmod 3$ erfordert $(p_3 \bmod 3, p_4 \bmod 3) \in \{(1,2), (2,1)\}$.
\end{proof}

\subsection{Kongruenzbedingungen modulo 4 und 8}

\begin{proposition}
\label{prop:mod4}
Falls $p_1 = 2$ und $p_2, p_3, p_4$ ungerade Primzahlen sind, so gilt für
jede ungerade Primzahl $p$ entweder $p \equiv 1 \pmod{4}$ oder
$p \equiv 3 \pmod{4}$.  Die Bedingung $p_2 \mid 2 p_3 p_4 - 1$ legt
zusätzliche Kongruenzen für $p_2, p_3, p_4$ modulo~$4$ fest, die im
Einzelfall zu untersuchen sind.
\end{proposition}

% -------------------------------------------------------
\section{Computergestützte Verifikation}
\label{sec:rechner}
% -------------------------------------------------------

\subsection{Suchalgorithmus}

Das Python-Modul \texttt{giuga\_4prim.py} des Specialist Mathematics Projects
implementiert eine systematische Suche nach Vier-Prim-Giuga-Pseudoprimzahlen.
Der Suchalgorithmus verwendet folgende Strategie:

\begin{enumerate}
  \item Für jede Primzahl $p_1$ bis zu einem Schwellenwert:
  \item Für jede Primzahl $p_2 > p_1$:
  \item Für jede Primzahl $p_3 > p_2$:
  \item Berechne die Menge der Kandidaten $p_4$: $p_4$ muss ein Primteiler
        von $p_1 p_2 p_3 - 1$ und größer als $p_3$ sein.
  \item Für jeden solchen Kandidaten $p_4$: überprüfe die Bedingungen
        \eqref{eq:G1}, \eqref{eq:G2}, \eqref{eq:G3}.
\end{enumerate}

Diese Vorgehensweise ist effizienter als eine naive Suche, da
Bedingung~\eqref{eq:G4} die Kandidaten für $p_4$ stark einschränkt.

\subsection{Ergebnisse}

\begin{theorem}[Computergestütztes Nicht-Existenzresultat]
\label{satz:rechner}
Die systematische Suche mit \texttt{giuga\_4prim.py} hat bestätigt, dass
kein Produkt $n = p_1 p_2 p_3 p_4 \le 10^{15}$ von vier verschiedenen
Primzahlen alle vier Giuga-Bedingungen gleichzeitig erfüllt.
\end{theorem}

\begin{proof}
Durch computergestützte Verifikation.  Alle Fälle mit $p_4 \le p_1 p_2 p_3 - 1$
und $n \le 10^{15}$ wurden mit exakter ganzzahliger Arithmetik unter Verwendung
der \texttt{sympy}-Bibliothek in Python~3.13 überprüft.
\end{proof}

\begin{remark}
Das Resultat in Satz~\ref{satz:rechner} ist konsistent mit, aber weit schwächer
als der theoretische untere Grenzwert von Borwein et al.\ \cite{Borwein1996},
der zeigt, dass jede Giuga-Pseudoprimzahl mehr als $10^{19907}$ Dezimalstellen
haben muss.  Unsere Berechnung liefert dennoch eine unabhängige Verifikation
und ein konkretes Verständnis der Dichte von Beinahe-Treffern.
\end{remark}

\subsection{Giuga-Zahlen als Beinahe-Treffer}

\begin{example}
\label{bsp:1722}
Betrachte $n = 1722 = 2 \cdot 3 \cdot 7 \cdot 41$ mit
$p_1 = 2 < p_2 = 3 < p_3 = 7 < p_4 = 41$.  Wir überprüfen:
\begin{itemize}
  \item $p_1 = 2$: $2 \mid 3 \cdot 7 \cdot 41 - 1 = 861 - 1 = 860$. \checkmark
  \item $p_2 = 3$: $3 \mid 2 \cdot 7 \cdot 41 - 1 = 574 - 1 = 573 = 191 \cdot 3$. \checkmark
  \item $p_3 = 7$: $7 \mid 2 \cdot 3 \cdot 41 - 1 = 246 - 1 = 245 = 35 \cdot 7$. \checkmark
  \item $p_4 = 41$: $41 \mid 2 \cdot 3 \cdot 7 - 1 = 42 - 1 = 41 = 1 \cdot 41$. \checkmark
\end{itemize}
Alle vier Giuga-Bedingungen sind erfüllt.  Also ist $1722$ eine Giuga-Zahl.
Ob $1722$ auch eine Giuga-Pseudoprimzahl ist, hängt davon ab, ob die
Summenkongruenz $\sum_{k=1}^{1721} k^{1721} \equiv -1 \pmod{1722}$ gilt.
\end{example}

\begin{remark}
Beispiel~\ref{bsp:1722} zeigt, dass es tatsächlich zusammengesetzte
Vier-Prim-Produkte gibt, die alle vier schwachen Giuga-Bedingungen erfüllen.
Die Pseudoprimzahl-Frage stellt die strengere Forderung der Summenkongruenz.
\end{remark}

\begin{example}
\label{bsp:keine}
Betrachte $p_1 = 2$, $p_2 = 3$, $p_3 = 11$, $p_4 = 23$:
\begin{itemize}
  \item $2 \mid 3 \cdot 11 \cdot 23 - 1 = 759 - 1 = 758$. \checkmark
  \item $3 \mid 2 \cdot 11 \cdot 23 - 1 = 506 - 1 = 505$?  $505 = 168 \cdot 3 + 1$. \texttimes
\end{itemize}
Dieser Kandidat scheitert bereits an der zweiten Bedingung.
\end{example}

% -------------------------------------------------------
\section{Offene Probleme}
\label{sec:offen}
% -------------------------------------------------------

Nach den Untersuchungen in dieser Arbeit verbleiben folgende Probleme offen.

\begin{conjecture}[Giugas Vermutung, {\cite{Giuga1950}}]
\label{verm:giuga}
Es gibt keine Giuga-Pseudoprimzahl.  Äquivalent: $n > 1$ ist prim genau dann,
wenn $\sum_{k=1}^{n-1} k^{n-1} \equiv -1 \pmod{n}$.
\end{conjecture}

\begin{conjecture}[Vier-Prim-Fall]
\label{verm:vierprim}
Es gibt keine Giuga-Pseudoprimzahl der Form $n = p_1 p_2 p_3 p_4$ mit vier
verschiedenen Primfaktoren.
\end{conjecture}

\begin{remark}
Vermutung~\ref{verm:vierprim} ist strikt schwächer als
Vermutung~\ref{verm:giuga}: Ein Beweis des Vier-Prim-Falls würde die bekannte
untere Schranke von drei Primfaktoren um einen Schritt erweitern, die
Gesamtvermutung aber nicht lösen.
\end{remark}

Weitere offene Probleme:

\begin{enumerate}
  \item \textbf{Eliminierung spezifischer Restklassen}: Lassen sich die
        Giuga-Bedingungen modulo kleiner Primzahlen für Vier-Prim-Produkte in
        bestimmten Kongruenzklassen als inkonsistent nachweisen?
  \item \textbf{Dichtefragen}: Was ist die Dichte von Vier-Prim-Produkten,
        die zwei oder drei der vier Giuga-Bedingungen erfüllen, und tendiert
        sie gegen null?
  \item \textbf{Zusammenhang mit Giuga-Zahlen}: Die bekannten Giuga-Zahlen
        $30, 858, 1722, \ldots$ sind Produkte, die die schwache Bedingung
        erfüllen.  Gibt es ein Muster, das sie mit Beinahe-Treffern für
        Pseudoprimzahlen verbindet?
  \item \textbf{Automatisierter Beweissuchlauf}: Kann der Vier-Prim-Fall durch
        einen vollautomatischen Beweiser auf Basis der in dieser Arbeit
        hergeleiteten modularen Arithmetikbedingungen gelöst werden?
  \item \textbf{$p$-adischer Zugang}: Können $p$-adische analytische Methoden,
        analog zur Iwasawa-Theorie elliptischer Kurven, Licht auf die Struktur
        der Lösungen des Vier-Prim-Giuga-Systems werfen?
\end{enumerate}

% -------------------------------------------------------
\section{Zusammenfassung}
\label{sec:zusammenfassung}
% -------------------------------------------------------

Wir haben eine systematische Untersuchung des Vier-Prim-Falls von Giugas
Vermutung aus dem Jahr 1950 durchgeführt.  Unsere Hauptresultate sind:

\begin{itemize}
  \item Vollständige Herleitung des Vier-Prim-Giuga-Gleichungssystems
        \eqref{eq:G1}--\eqref{eq:G4}.
  \item Die Totient-Identität $\varphi(n) = n - e_3 + e_2 - e_1 + 1$
        (Satz~\ref{satz:totient}), die die algebraische Struktur des Problems
        beleuchtet.
  \item Mehrere notwendige arithmetische Einschränkungen, darunter die Schranke
        $p_4^2 \le p_1 p_2 p_3$ (Satz~\ref{satz:quadratschranke}).
  \item Paritäts- und Kongruenzanalysen, einschließlich der automatischen
        Erfüllung der geraden-Faktor-Bedingung (Satz~\ref{satz:p1auto}).
  \item Computergestützte Nicht-Existenz für $n \le 10^{15}$
        (Satz~\ref{satz:rechner}).
\end{itemize}

Der Vier-Prim-Fall von Giugas Vermutung---und Giugas Vermutung selbst---
\emph{bleibt offen}.  Wir hoffen, dass die strukturelle Analyse dieser Arbeit
zu künftigen Fortschritten beiträgt.

\section*{Danksagung}

Der Autor dankt dem Specialist Mathematics Project (Build~131) für die
bereitgestellte Recheninfrastruktur und das mathematische Rahmenwerk, das
diese Untersuchung ermöglicht hat.

% -------------------------------------------------------
\begin{thebibliography}{10}

\bibitem{Giuga1950}
G.~Giuga,
\emph{Su una presumibile propriet\`a caratteristica dei numeri primi},
Ist.\ Lombardo Sci.\ Lett.\ Rend.\ Cl.\ Sci.\ Mat.\ Nat.\ \textbf{83}
(1950), 511--528.

\bibitem{Borwein1996}
J.~M.~Borwein, P.~B.~Borwein, R.~Girgensohn und S.~Parnes,
\emph{Giuga's conjecture on primality},
Amer.\ Math.\ Monthly \textbf{103} (1996), Nr.\ 1, 40--50.

\bibitem{Agoh1995}
T.~Agoh,
\emph{On Giuga's conjecture},
Manuscripta Math.\ \textbf{87} (1995), 501--510.

\bibitem{Tipu2006}
T.~Tipu,
\emph{A note on Giuga's conjecture},
Canad.\ Math.\ Bull.\ \textbf{49} (2006), Nr.\ 3, 472--480.

\bibitem{Bednarek2014}
M.~Bednarek,
\emph{On the size of Giuga pseudoprimes},
Preprint (2014).

\bibitem{FuhrmannPaper1}
M.~Fuhrmann,
\emph{Keine zusammengesetzte Zahl mit genau drei Primfaktoren ist eine
Giuga-Pseudoprimzahl: Ein vollständiger elementarer Beweis},
Specialist Mathematics Project, Build~43 (2025).

\bibitem{FuhrmannPaper4}
M.~Fuhrmann,
\emph{Giuga-Zahlen und Quadratfreiheit},
Specialist Mathematics Project, Build~50 (2025).

\bibitem{FuhrmannPaper5}
M.~Fuhrmann,
\emph{Kein Semiprimes ist eine Giuga-Zahl},
Specialist Mathematics Project, Build~50 (2025).

\bibitem{Hardy1979}
G.~H.~Hardy und E.~M.~Wright,
\emph{An Introduction to the Theory of Numbers},
fünfte Auflage, Oxford University Press, Oxford, 1979.

\bibitem{Ribenboim1996}
P.~Ribenboim,
\emph{The New Book of Prime Number Records},
Springer, New York, 1996.

\end{thebibliography}

\end{document}
